%% article.tex -- abstract and section plan for an article announcing Odum BASIC
%% Simulations, for submission to Ecological Modelling.
%%
%% STATUS: plan. The abstract and highlights are draft prose. Everything in a
%% "plan" box says what a section must argue and what evidence carries it; it is
%% not finished text. Items marked [VERIFY] are facts or references not yet
%% checked against a primary source. Nothing marked [VERIFY] goes into the
%% submitted manuscript unchecked.
%%
%% Build: pdflatex article.tex (twice, for cross-references), from docs/.
%% References are inline (thebibliography) in Elsevier's author-year style, as in
%% the companion energese draft, so no BibTeX run is needed.
%%
%% Submission format: Elsevier wants elsarticle.cls. This drafts in `article` and
%% switches at submission:
%%   \documentclass[preprint,12pt,authoryear]{elsarticle}
%% with \begin{frontmatter}...\end{frontmatter} replacing \maketitle, and the
%% highlights moved to their own file.
%%
%% Paragraphs are single long lines: do not hard-wrap prose in this file.

\documentclass[11pt,a4paper]{article}
\usepackage[margin=1in]{geometry}
\usepackage[T1]{fontenc}
\usepackage{lmodern}
\usepackage{booktabs}
\usepackage{tabularx}
\usepackage{listings}
\usepackage{xcolor}
\usepackage{framed}
\usepackage[round]{natbib}
\usepackage{hyperref}

\lstset{basicstyle=\ttfamily\footnotesize, columns=fullflexible, breaklines=true, frame=single, framesep=4pt, xleftmargin=6pt}

% A plan block is what a section must argue, not its text. A left bar rather than
% a shaded box, because a box cannot break across pages and some plans are long.
\newenvironment{plan}{\begin{leftbar}\small\setlength{\parskip}{0.45em}\noindent\textbf{Plan.}\ }{\end{leftbar}}
\newcommand{\target}[1]{\par\noindent\emph{Target length: #1 words.}}
\newcommand{\verify}{\textbf{[VERIFY]}}

\title{Odum BASIC Simulations: a provenance-tracked archive in which the published simulation listings of systems ecology run again}
\author{S. Maud\thanks{[Affiliation and ORCID to add.]}}
\date{Plan --- September 2026}

\begin{document}
\maketitle

%% ----------------------------------------------------------------------------
\section*{About this plan}

This document is the abstract and the section-by-section plan for a paper announcing the project to \emph{Ecological Modelling}. It fixes what each section argues, what evidence carries the argument, and how long it should be, so that drafting can proceed section by section without the paper drifting in length or scope. Two things make the paper unusual for the journal and shape the plan throughout: the software was built largely by a large language model (LLM) coding agent, which Elsevier requires to be described as method rather than hidden \citep{elsevierAI}; and one of the contribution routes it describes hands transcription of new programs to a second agent. Both are presented as method, with their controls and their failures, not as novelty.

\paragraph{Target journal and type.} \emph{Ecological Modelling} published the memorial issue on H.T. Odum's contributions \citep{brown2004}, and the companion \texttt{energese} paper on typesetting Energy Systems Language diagrams is being prepared for the same journal. The two should cite each other and must not overlap: \texttt{energese} is about drawing the diagram, this paper is about running the model. Article type: Research Paper unless the current guide for authors offers a software or short-communication type that fits better \verify. The guide for authors could not be retrieved while writing this plan; check the article types, the abstract limit and the highlights rules before submission \verify.

\paragraph{Length.} About 6{,}500 words of body text, three figures, five tables and one listing. That is long enough to cover the archive, the software and the LLM method properly, and short enough that the paper does not become a user manual; anything that reads like documentation belongs in the repository, which the paper cites.

\begin{table}[h]
\centering\small
\begin{tabular}{@{}llr@{}}
\toprule
& Section & Words \\
\midrule
1 & Introduction & 700 \\
2 & Background: simulation in Odum's systems ecology & 800 \\
3 & Design principles & 600 \\
4 & Implementation & 1{,}000 \\
5 & Contributing programs & 900 \\
6 & Building the software with an LLM coding agent & 1{,}000 \\
7 & Validation & 500 \\
8 & Discussion & 750 \\
9 & Conclusions & 250 \\
\midrule
& Total (excluding abstract, declarations, references) & 6{,}500 \\
\bottomrule
\end{tabular}
\caption{Length budget.}
\end{table}

\paragraph{Before this can be submitted.} The software exists and is deployed, but the paper would overclaim if submitted today. Table~\ref{tab:prereq} lists what must be true first. The first item matters most: at present the archive holds five programs written for the repository as worked examples and no transcription of a published listing, and an archive paper with nothing archived is not yet a paper.

\begin{table}[h]
\centering\small
\begin{tabularx}{\linewidth}{@{}Xl@{}}
\toprule
Prerequisite & Status (18 Sept 2026) \\
\midrule
A substantial set of \texttt{verbatim} transcriptions from Odum's books, each with its citation, and a rights decision for every reproduced diagram & 0 transcribed \\
A licence file in the repository (without one the code is not legally reusable, whatever the paper says) & missing \\
A citable release: \texttt{CITATION.cff} and a Zenodo DOI, as for \texttt{energese} & missing \\
The GW-BASIC/QuickBASIC statements the later listings use, especially \texttt{PSET}/\texttt{LINE} as plot data, so that the graphical minimodels run unaltered & planned \\
The \texttt{jules} route configured: the Jules GitHub app installed, and a section in \texttt{AGENTS.md} for cloud agents, which cannot use the Apple \texttt{container} toolchain the file currently prescribes & not configured \\
Whether GitHub prefills issue-form dropdowns from a URL, checked with a real submission & unverified \\
The one intermittently failing end-to-end test diagnosed & open \\
Current guide-for-authors requirements checked (type, abstract limit, highlights) & \verify \\
\bottomrule
\end{tabularx}
\caption{What must be true before submission.}
\label{tab:prereq}
\end{table}

%% ----------------------------------------------------------------------------
\begin{abstract}
H.T. Odum and his colleagues published many of their ecological simulation models as BASIC listings: short numbered programs, printed beside an Energy Systems Language diagram, meant to be typed in and run. The machines they were written for are gone, and the listings survive only on paper, so the models can be read but are rarely run, checked or extended. We present Odum BASIC Simulations, an open archive and browser workbench in which these listings run unaltered. A BASIC interpreter written for the purpose runs in the browser; the table a program prints is recovered and plotted, and the printed output stays visible beside the plot. Every program is stored with a machine-checked provenance record: a citation in BibLaTeX form, a declared fidelity to the published listing (verbatim, corrected, adapted or original) and, where available, the program's diagram with the rights under which it is reproduced. The build refuses any program whose record is incomplete. Programs can be contributed as a pull request, through a GitHub issue form processed by an automated workflow, from an in-browser workspace that fills in that form, or by labelling an issue for a coding agent, which transcribes the source into archive files for human review. The software itself was built largely by an LLM coding agent working under human direction, test-first, against written repository standards, with every change reviewed and merged by a person; we describe that process, its controls and the failures they caught. For the first mini-model the interpreter reproduces the closed-form Euler solution to print precision (relative difference $2\times10^{-9}$). The archive is openly available at \url{https://energese-project.github.io/Odum-basic-simulations/}.
\end{abstract}

\begin{plan}
About 250 words as drafted, which should sit inside the journal's limit \verify. The abstract states the problem (listings exist only on paper), the contribution (runs unaltered, provenance enforced, four contribution routes), the method disclosure (LLM-built, with controls), one quantitative result, and availability. Once transcriptions exist, replace the validation sentence with the stronger result from Section~\ref{sec:validation}: agreement with output printed in the source.
\end{plan}

\paragraph{Highlights} (five, each at most 85 characters including spaces):
\begin{itemize}
\item Odum's BASIC simulation listings run unaltered in a web browser
\item Every program carries a machine-checked citation and fidelity record
\item Published diagrams are archived only with their reproduction rights stated
\item Programs can be contributed from the browser, a form, or a coding agent
\item The software was built by an LLM agent under test-first human review
\end{itemize}

\paragraph{Keywords:} systems ecology; H.T. Odum; energy systems language; simulation; BASIC; software preservation; research software; large language models

%% ----------------------------------------------------------------------------
\section{Introduction}
\label{sec:intro}

\begin{plan}
Three moves, in this order.

\emph{The asset.} Odum's teaching and research books carry dozens of complete simulation models as BASIC listings, each paired with its Energy Systems Language diagram \citep{odum1983,odumOdum2000,odumMinimodels} \verify{} (count them per book before stating a number). They are among the most concrete forms in which systems ecology was published: the model is not described, it is given, and it can be executed.

\emph{The loss.} A listing on paper is a model that cannot be run. The interpreters are gone, the dialects have drifted, retyping introduces errors nobody detects, and a model that cannot be run cannot be checked, compared or reused. This is the general problem of software collapse \citep{hinsen2019} in its purest form: the code survived and everything it depended on did not.

\emph{The contribution.} (i) An interpreter that runs the published listings as printed, with plotting recovered from their text output. (ii) An archive in which each program's provenance is a checked record rather than a comment, including how faithful it is to the page and on what basis its diagram is reproduced. (iii) Contribution routes that let someone with a book and no programming background add a program. (iv) A documented, auditable account of building such software with an LLM coding agent, including how an agent can also serve as a transcription route.

End with a one-paragraph map of the paper.
\end{plan}
\target{700}

%% ----------------------------------------------------------------------------
\section{Background: simulation in Odum's systems ecology}
\label{sec:background}

\begin{plan}
\emph{From circuits to code.} Odum's energy circuit language began as analogue electrical circuits \citep{odum1960} and became a notation for systems of any kind \citep{odum1972,odum1983,odum1994}. Simulation was always part of it: a diagram specifies a set of equations, and the equations were run, first on analogue computers and later as BASIC on microcomputers. Keep this short and cite the \texttt{energese} companion for the language itself.

\emph{BASIC as the teaching medium.} Why BASIC \citep{kemenyKurtz1964}: it was on every microcomputer students could reach, and a listing could be printed in a book and typed in. Describe the shape of an Odum listing (parameters as \texttt{LET}, an Euler loop over \texttt{DT}, a \texttt{PRINT} per step, sometimes screen graphics) using a short real example. Identify the dialects in the sources: Commodore-era BASIC V2 in the earlier material, GW-BASIC/QuickBASIC with \texttt{SCREEN}/\texttt{PSET}/\texttt{LINE} in the later \verify{} (confirm per book).

\emph{What exists today, and why it is not enough.} Emulators run the old machines but need ROM images that cannot be redistributed, and they run the machine, not the model; general system-dynamics tools (Stella, Vensim, Insight Maker) can re-express a model but not run the listing, so the published artefact itself is never checked \verify{} (cite each tool properly). Prior re-implementations of individual Odum models \verify{} (search the literature; cite any that exist). The gap: no archive in which the listings themselves run, with their provenance.

\emph{Provenance and research software.} Place the design against FAIR for research software \citep{barker2022} and software citation principles \citep{smith2016}. The archive adds a dimension those frameworks do not name: fidelity, meaning how close an executable artefact is to the published one.
\end{plan}
\target{800}

%% ----------------------------------------------------------------------------
\section{Design principles}
\label{sec:design}

\begin{plan}
Four principles, each stated with the design decision it forced. This section is the paper's argument about \emph{how} such an archive should be built, and should stay short and general; Section~\ref{sec:implementation} is the how-it-was-done.

\emph{The listing is the product.} When a published listing and the interpreter disagree, the interpreter changes. Consequences: no rewriting listings into a modern language; plotting recovered from printed text rather than by requiring an output convention; unsupported statements treated as interpreter work, not as edits to the source.

\emph{Provenance is data, and incomplete provenance is a build failure.} A listing whose origin cannot be read looks authoritative and is not, so it is refused. Introduce \emph{fidelity} (Table~\ref{tab:fidelity}). It has no default, because it changes what every other field means, and the one rule no automated route may break is that a person chooses it.

\emph{Built to outlive its authors.} A static site with no server or database; the browser platform first; two runtime dependencies (Monaco, Chart.js), bundled rather than loaded from a CDN; build tooling kept out of the shipped artefact. Argue why: the archive's value is realised over decades, long after any maintainer has moved on \citep{hinsen2019}.

\emph{Test-first, from the outside.} Interpreter tests run a listing and assert on what it printed; archive tests are phrased as the mistakes a contributor could make. That makes the tests survive a rewrite of the implementation, and it matters again in Section~\ref{sec:llm}, where the tests are the main check on an implementation the human author did not write line by line.
\end{plan}

\begin{table}[h]
\centering\small
\begin{tabularx}{\linewidth}{@{}lXl@{}}
\toprule
Fidelity & Meaning & Also requires \\
\midrule
\texttt{verbatim} & Transcribed from the source character for character, including its errors & a source \\
\texttt{corrected} & Transcribed, with misprints fixed & a source; every change listed \\
\texttt{adapted} & The published model, rewritten to run here & a source \\
\texttt{original} & Written for the archive; not a published listing & no source \\
\bottomrule
\end{tabularx}
\caption{Fidelity levels. The contributor chooses the lowest level that is true; a broken published listing is archived as \texttt{verbatim} and its fix as a separate \texttt{corrected} entry, so both are on record.}
\label{tab:fidelity}
\end{table}
\target{600}

%% ----------------------------------------------------------------------------
\section{Implementation}
\label{sec:implementation}

\begin{plan}
Describe what exists, at the level a modeller needs in order to trust and use it. Point to the repository for everything else.

\emph{The interpreter} (about 800 lines of TypeScript). Statements and functions supported (Commodore BASIC V2: \texttt{PRINT}, \texttt{LET}, \texttt{FOR}/\texttt{NEXT}, \texttt{IF}/\texttt{THEN}, \texttt{GOSUB}, \texttt{INPUT}, \texttt{DIM}, \texttt{DATA}/\texttt{READ}, maths including \texttt{EXP} and \texttt{LOG}), and those planned. It runs in a Web Worker so a runaway program can be stopped; it knows nothing of the page, which is why the same code runs in the browser and under the unit tests.

\emph{Recovering the plot from printed text.} State the heuristic exactly: a run of wholly numeric lines of equal width is a table, the first column is the independent variable, and the non-numeric line above supplies the names. Say why it is a heuristic on purpose: requiring an output convention would mean rewriting the published programs. One y axis, never two, and the printed table always stays visible beside the chart, so the plot is never the only reading of a run.

\emph{The workbench} (Figure~\ref{fig:workbench}): explorer, editor, diagram, plot and output. Keep this to one paragraph; it is the figure's caption, not a manual.

\emph{The archive.} One directory, three files per program (\texttt{.bas}, \texttt{.json}, optional diagram), no manifest. At build time every sidecar is validated (Listing~\ref{lst:sidecar}); the raw files are published at stable URLs so that a citation can point at the listing itself, and the tests check that the published file is byte-identical to the one the app runs.

\emph{Diagrams and rights.} The diagram is the other half of what Odum published, and usually the easier half to check a model against. It is archived only with its rights recorded (Table~\ref{tab:rights}); a reproduced figure must cite its source; the build rejects an image nobody has declared, a mismatch between an image's content and its file type, and anything over 2 MB. SVG is refused because an SVG opened from the site can run scripts. Keep this one paragraph; the rights argument is in Section~\ref{sec:discussion}.

Figures: \ref{fig:workbench} (workbench), \ref{fig:minimodel} (diagram and run of the first mini-model). Tables: \ref{tab:rights}.
\end{plan}

\begin{lstlisting}[caption={The sidecar for the first mini-model, \texttt{programs/charge-discharge.json}, with the notes shortened and $K_1\cdot Q$ written as \texttt{K1*Q}. The build fails if any required field is missing or inconsistent.},label={lst:sidecar}]
{
  "title": "Charge And Discharge",
  "description": "One tank filled at a constant rate and drained in proportion to storage.",
  "fidelity": "original",
  "notes": "The model form is the first mini-model in the systems-ecology sequence. ...",
  "tags": ["mini-model", "storage", "steady-state"],
  "diagram": {
    "file": "charge-discharge.png",
    "caption": "Energy systems diagram: a source J fills the storage Q, which drains through K1*Q to a heat sink.",
    "rights": {
      "basis": "own-work",
      "statement": "Drawn for this repository in Odum's energy systems symbols. Not reproduced from a publication."
    }
  }
}
\end{lstlisting}

\begin{table}[h]
\centering\small
\begin{tabularx}{\linewidth}{@{}lXl@{}}
\toprule
Rights basis & Meaning & Also requires \\
\midrule
\texttt{own-work} & Drawn for the archive & --- \\
\texttt{public-domain} & Out of copyright, or never in it & a source \\
\texttt{licensed} & Under a licence that permits it, named in the statement & a source \\
\texttt{permission} & The rights holder agreed; who and when are recorded & a source \\
\texttt{fair-use} & Fair use or fair dealing: reproduced for scholarship or review beside the listing it documents & a source \\
\bottomrule
\end{tabularx}
\caption{The basis on which a diagram is published. Like fidelity it has no default, and every basis also requires a written statement.}
\label{tab:rights}
\end{table}

\begin{figure}[h]
\centering
\fbox{\parbox{0.9\linewidth}{\centering\vspace{2em}[Screenshot of the workbench: explorer with Archive and My programs, the listing in the editor, the diagram above the plot, and the printed output. Generate it from the end-to-end suite, as \texttt{energese} generates its figures from its test harness, so that the figure is reproducible from the repository at a tagged commit.]\vspace{2em}}}
\caption{The workbench.}
\label{fig:workbench}
\end{figure}

\begin{figure}[h]
\centering
\fbox{\parbox{0.9\linewidth}{\centering\vspace{2em}[Two panels: (a) the Energy Systems Language diagram of the charge-and-discharge mini-model, \texttt{programs/charge-discharge.png}; (b) the storage $Q$ and outflow $K_1Q$ plotted from the program's printed table. Once transcriptions exist, replace this with a published minimodel and its original figure beside the recovered plot.]\vspace{2em}}}
\caption{A model as published: its diagram, and the run recovered from its listing.}
\label{fig:minimodel}
\end{figure}
\target{1{,}000}

%% ----------------------------------------------------------------------------
\section{Contributing programs}
\label{sec:contributing}

\begin{plan}
The archive grows only if people with the books can add to it, and most of them are ecologists, not programmers. Present four routes in order of decreasing technical demand (Table~\ref{tab:routes}, Figure~\ref{fig:routes}). The single point to make: every route ends in the same validator and the same human review, so lowering the barrier does not lower the bar.

\emph{A pull request} with the three files, for contributors who use git.

\emph{A GitHub issue form.} The contributor pastes the listing and types the citation; a workflow parses the issue, builds the files, checks the whole archive with the build's own validator, and either comments every problem at once or opens a pull request crediting the contributor as co-author. Editing the issue re-runs it. Note the two security properties, briefly: the issue text is never passed to a shell, and a diagram is fetched only from GitHub's own attachment hosts.

\emph{The in-browser workspace.} A contributor writes the program in the workbench itself, in a private area of the browser's storage (the Origin Private File System), runs it, and fills in the same fields, checked as they type. \emph{Submit} opens the issue form already filled in. The archive itself is never copied into the workspace: a copy would go stale when an entry is corrected upstream, and an edited copy of a \texttt{verbatim} listing is not verbatim, so copying an archive program clears its fidelity and records where it came from. There is no sign-in, because a static site cannot complete GitHub's OAuth exchange without a server holding the client secret \citep{githubOAuth}; say this in one sentence as a deliberate trade.

\emph{Labelling an issue for a coding agent.} A contributor opens an issue with the source: a scan of the page, or a reference and a pasted listing. A maintainer adds the \texttt{jules} label, and Jules \citep{jules}, Google's asynchronous coding agent, picks the issue up, reads the repository's standards in \texttt{AGENTS.md}, transcribes the listing and metadata into the three files, runs the tests, and opens a pull request. State the limits precisely. The agent does the clerical work; it must not choose the fidelity, which comes from the contributor or is confirmed by the reviewer against the page photograph; and a transcription from a scan is exactly where an undetected error enters an archive, so the reviewer's check against the page is mandatory. (Before submission: configure the route and report its first real uses, including any errors the reviewer caught; see Table~\ref{tab:prereq}.)
\end{plan}

\begin{table}[h]
\centering\small
\begin{tabularx}{\linewidth}{@{}lXXX@{}}
\toprule
Route & Contributor needs & Automated & Human step \\
\midrule
Pull request & git, the file formats & build validation, full test suite & review, merge \\
Issue form & a GitHub account & files built and validated; PR opened; CI run & review, merge \\
Workspace & a browser (GitHub account to submit) & live validation; form prefilled & drag in the diagram; review, merge \\
\texttt{jules} label & a GitHub account & transcription, files, tests, PR by the agent & maintainer applies the label; fidelity confirmed against the page; review, merge \\
\bottomrule
\end{tabularx}
\caption{Contribution routes. All four end in the same validation and the same review.}
\label{tab:routes}
\end{table}

\begin{figure}[h]
\centering
\fbox{\parbox{0.9\linewidth}{\centering\vspace{2em}[Flow diagram: four entry points (pull request, issue form, workspace, \texttt{jules} label) converging on one validator (the build's \texttt{parseProgramMeta} and \texttt{readCatalog}), then CI, then human review, then the published archive. Draw in TikZ.]\vspace{2em}}}
\caption{Four ways in, one validator, one review.}
\label{fig:routes}
\end{figure}
\target{900}

%% ----------------------------------------------------------------------------
\section{Building the software with an LLM coding agent}
\label{sec:llm}

\begin{plan}
This is a methods section, written to Elsevier's requirement that AI used in the research process be described in the methods with the tool, version and developer \citep{elsevierAI}. The tone is an account of a method, with its controls and its failures, not an advertisement. It should let a reader judge how far to trust code they did not see a person write.

\emph{Tools.} Claude Code \citep{claudeCode}, Anthropic's agentic coding tool, running the model Claude Opus~5 (\texttt{claude-opus-5}), in September 2026, from the author's editor. Jules is used only as a contribution route (Section~\ref{sec:contributing}), not to build the software.

\emph{Division of labour.} The author set the requirements and made every design decision; the agent proposed designs, wrote the code, tests and documentation, ran the checks and opened the pull requests. Give real decisions that stayed with the author: publishing scanned diagrams only with rights recorded; submitting through an issue form rather than an OAuth service; keeping the archive out of the workspace; the interface layout, revised from a marked-up screenshot. The agent was not permitted to merge its own work, and every change reached the main branch through a pull request the author reviewed and merged.

\emph{Controls.} Written repository standards in \texttt{AGENTS.md} (architecture, test-first, dependency policy, security rules), the same file Jules reads, so the two agents work to one standard; tests written before code; continuous integration running the full suite against both the development server and the production build; one isolated git worktree per task; and a co-author trailer on every agent-assisted commit, so the history records which changes the agent took part in (Table~\ref{tab:llm}).

\emph{What the controls caught.} The most useful evidence is failures, each drawn from the repository history. (a) A style rule leaking between interface components, which the tests did not catch and screenshots did. (b) The issue-form parser attaching an unrelated section's text to the diagram rights statement, found by a dry run against a synthetic issue before deployment. (c) A workflow expression using a context that is unavailable where it was used, found by a workflow linter. (d) An intermittently failing end-to-end test the agent could not diagnose, which it reported in the pull request instead of suppressing. (e) A behaviour of GitHub it could not verify (prefilling dropdowns from a URL), which it flagged as unverified instead of asserting. The pattern to draw out: the agent was reliable where a check existed and candid where one did not, which puts the burden on designing checks.

\emph{Limits and threats.} Git trailers attribute commits, not lines, so the agent's share of the code is known only coarsely. The process cannot be reproduced exactly, because the model is not deterministic and will be superseded; what is reproducible is the artefact, its tests and its history. Review is the bottleneck, and a single reviewer is a single point of failure. The interpreter is the component where an undetected error does the most harm, which is why it carries the analytical check in Section~\ref{sec:validation} and why agreement with printed source output is the validation that matters most. Place this against the software-engineering evidence on LLM coding agents \citep{chen2021,jimenez2024}.
\end{plan}

\begin{table}[h]
\centering\small
\begin{tabularx}{\linewidth}{@{}Xr@{}}
\toprule
Quantity (main branch at commit \texttt{4edbbdd}, 18 Sept 2026) & Value \\
\midrule
Pull requests merged & 8 \\
Non-merge commits & 9 \\
\quad of which carry a Claude co-author trailer & 8 \\
\quad the remaining one & GitHub's initial commit \\
TypeScript/JavaScript source lines (excluding tests, configuration and the vendored framework) & 3{,}610 \\
\quad of which the BASIC interpreter & 798 \\
CSS and HTML lines & 1{,}374 \\
Test lines & 1{,}608 \\
Unit tests & 136 \\
End-to-end tests (41, each run against the development and the production build) & 82 \\
Elapsed time from first to latest merge & about 12 hours \\
\bottomrule
\end{tabularx}
\caption{The development record. Recompute at the tagged release the paper cites.}
\label{tab:llm}
\end{table}
\target{1{,}000}

%% ----------------------------------------------------------------------------
\section{Validation}
\label{sec:validation}

\begin{plan}
Two kinds of evidence, the second decisive.

\emph{Against analytical results.} The charge-and-discharge mini-model, $\mathrm{d}Q/\mathrm{d}t = J - K_1Q$, discretised by the listing's own Euler step, has the closed form $Q_n = (J/K_1)\,[1-(1-K_1\Delta t)^n]$. The interpreter's printed output matches this recurrence to print precision, so the difference from the exact solution $Q(t) = (J/K_1)(1-e^{-K_1t})$ is the listing's method and not the interpreter's (Table~\ref{tab:validation}). Do the same for logistic growth and for the two tanks in series.

\emph{Against output printed in the sources.} Where a book prints a program's output, as a table or a graph, the \texttt{verbatim} transcription must reproduce it. This is the result that validates the archive rather than the arithmetic, and it depends on the transcriptions listed in Table~\ref{tab:prereq}. Report agreement per program, and treat every disagreement as a finding: a misprint in the book, a transcription error, or an interpreter defect, each resolved on the record through the \texttt{corrected} mechanism.

\emph{The test suite} as supporting evidence: listing-level interpreter tests, archive-rule tests, and end-to-end tests against the production build, including that each published file is byte-identical to what the app runs.
\end{plan}

\begin{table}[h]
\centering\small
\begin{tabular}{@{}lr@{}}
\toprule
Charge and discharge: $J=100$, $K_1=0.1$, $\Delta t=0.5$, $0\le t\le60$ & \\
\midrule
Rows printed & 121 \\
Maximum relative difference, interpreter vs.\ Euler recurrence & $2.2\times10^{-9}$ \\
Maximum relative difference, interpreter vs.\ exact solution & $2.5\times10^{-2}$ \\
$Q(10)$: interpreter / recurrence / exact & 641.514 / 641.514 / 632.121 \\
\bottomrule
\end{tabular}
\caption{The interpreter reproduces the listing's arithmetic to print precision; the remaining difference from the exact solution is the listing's Euler step. Computed at commit \texttt{4edbbdd}; add this comparison to the test suite so the figure is regenerated rather than transcribed.}
\label{tab:validation}
\end{table}
\target{500}

%% ----------------------------------------------------------------------------
\section{Discussion}
\label{sec:discussion}

\begin{plan}
\emph{What running the listings makes possible.} Checking a published model against its own claims; teaching with the originals; comparing a listing with its diagram, which is also the bridge to \texttt{energese} and to GSSK, where one model file both draws and simulates. Keep speculation short.

\emph{Fidelity as a general idea.} Argue that recording how close an executable artefact is to the published one belongs in research-software practice beyond this archive, as a companion to FAIR \citep{wilkinson2016,barker2022}.

\emph{Figures, copyright and fair dealing.} The diagrams are someone's work. The design makes the rights decision explicit, per figure, recorded and reviewed, rather than implicit in whether someone dared to upload a scan. Say plainly what the archive does not decide for anyone, and that a basis can be revised if a rights holder objects. Consult a copyright adviser on the fair-use and fair-dealing wording before submission \verify.

\emph{Agents as contributors to scholarly archives.} Two roles, builder and transcriber, with different risks. The first is checked by tests; the second only by a person reading the page. The archive's rules (no default fidelity, mandatory review, provenance on every file) are what make an agent contribution acceptable, and they would be as necessary without agents.

\emph{Limitations.} Dialect coverage; the plotting heuristic; single-axis plots; no sign-in; dependence on GitHub for contribution (though not for use: the site is static and the archive is plain files); a single maintainer.
\end{plan}
\target{750}

%% ----------------------------------------------------------------------------
\section{Conclusions}
\label{sec:conclusions}

\begin{plan}
One paragraph of what exists and what it enables; one of what is needed from the community: transcriptions, rights decisions for figures, and reports of listings that do not yet run. End with the URL and the DOI.
\end{plan}
\target{250}

%% ----------------------------------------------------------------------------
\section*{Software and data availability}

\begin{plan}
Name: Odum BASIC Simulations. Developer: S. Maud. Repository: \url{https://github.com/energese-project/Odum-basic-simulations}. Live: \url{https://energese-project.github.io/Odum-basic-simulations/}. Licence: [to add; see Table~\ref{tab:prereq}]. Archived release and DOI: [Zenodo; to mint]. Language: TypeScript; runs in any current browser with no installation. Requirements to build: Node and the container image described in the repository. Size: see Table~\ref{tab:llm}. The archive's programs are plain files in \texttt{programs/}, each with its provenance record, citable individually by URL.
\end{plan}

\section*{CRediT authorship contribution statement}
\textbf{S. Maud:} Conceptualization, Methodology, Software, Validation, Writing -- original draft, Writing -- review \& editing. [Add co-authors and their roles.]

\section*{Declaration of competing interest}
[To complete. State whether any author has a relationship with Anthropic or Google.]

\section*{Declaration of generative AI and AI-assisted technologies in the manuscript preparation process}

\begin{plan}
Required by Elsevier as a separate section before the references, in their template wording \citep{elsevierAI}. This covers use in \emph{writing}; use in building the software is method, and is described in Section~\ref{sec:llm}. AI tools are not authors. Draft, to be revised to match what is actually used when the manuscript is written:
\end{plan}

During the preparation of this work the author used Claude Code (Anthropic; model Claude Opus~5) in order to draft the section plan and to assemble the development record reported in Table~\ref{tab:llm}. After using this tool, the author reviewed and edited the content as needed and takes full responsibility for the content of the published article.

\section*{Acknowledgements}
[To complete.]

%% ----------------------------------------------------------------------------
\begin{thebibliography}{99}

%% --- Odum and the Energy Systems Language -----------------------------------
\bibitem[Odum(1960)]{odum1960}
Odum, H.T., 1960. Ecological potential and analogue circuits for the ecosystem. \emph{American Scientist} 48(1), 1--8. %% [VERIFY] volume, pages

\bibitem[Odum(1972)]{odum1972}
Odum, H.T., 1972. An energy circuit language for ecological and social systems: its physical basis. In: Patten, B.C. (Ed.), \emph{Systems Analysis and Simulation in Ecology}, Vol.~2. Academic Press, New York, pp. 139--211. %% [VERIFY] pages

\bibitem[Odum(1983)]{odum1983}
Odum, H.T., 1983. \emph{Systems Ecology: An Introduction}. Wiley, New York.

\bibitem[Odum(1994)]{odum1994}
Odum, H.T., 1994. \emph{Ecological and General Systems: An Introduction to Systems Ecology}. University Press of Colorado, Niwot.

\bibitem[Odum and Odum(2000)]{odumOdum2000}
Odum, H.T., Odum, E.C., 2000. \emph{Modeling for All Scales: An Introduction to System Simulation}. Academic Press, San Diego.

\bibitem[Odum and Odum(1989)]{odumMinimodels}
Odum, H.T., Odum, E.C., 1989. \emph{Computer Minimodels and Simulation Exercises for Science and Social Science}. Center for Wetlands, University of Florida, Gainesville. %% [VERIFY] title, year, publisher -- from memory, not checked

\bibitem[Brown and Ulgiati(2004)]{brown2004}
Brown, M.T., Ulgiati, S., 2004. Energy quality, emergy, and transformity: H.T. Odum's contributions to quantifying and understanding systems. \emph{Ecological Modelling} 178(1--2), 201--213.

%% --- BASIC and software preservation ----------------------------------------
\bibitem[Kemeny and Kurtz(1964)]{kemenyKurtz1964}
Kemeny, J.G., Kurtz, T.E., 1964. \emph{BASIC: A Manual for BASIC, the Elementary Algebraic Language Designed for Use with the Dartmouth Time Sharing System}. Dartmouth College Computation Center, Hanover, NH. %% [VERIFY]

\bibitem[Hinsen(2019)]{hinsen2019}
Hinsen, K., 2019. Dealing with software collapse. \emph{Computing in Science \& Engineering} 21(3), 104--108. %% [VERIFY] pages

\bibitem[Wilkinson et~al.(2016)]{wilkinson2016}
Wilkinson, M.D., et al., 2016. The FAIR Guiding Principles for scientific data management and stewardship. \emph{Scientific Data} 3, 160018.

\bibitem[Barker et~al.(2022)]{barker2022}
Barker, M., Chue Hong, N.P., Katz, D.S., Lamprecht, A.-L., Martinez-Ortiz, C., Psomopoulos, F., Harrow, J., Castro, L.J., Gruenpeter, M., Martinez, P.A., Honeyman, T., 2022. Introducing the FAIR Principles for research software. \emph{Scientific Data} 9, 622.

\bibitem[Smith et~al.(2016)]{smith2016}
Smith, A.M., Katz, D.S., Niemeyer, K.E., FORCE11 Software Citation Working Group, 2016. Software citation principles. \emph{PeerJ Computer Science} 2, e86.

%% --- LLMs and coding agents -------------------------------------------------
\bibitem[Chen et~al.(2021)]{chen2021}
Chen, M., et al., 2021. Evaluating large language models trained on code. arXiv:2107.03374.

\bibitem[Jimenez et~al.(2024)]{jimenez2024}
Jimenez, C.E., Yang, J., Wettig, A., Yao, S., Pei, K., Press, O., Narasimhan, K., 2024. SWE-bench: Can language models resolve real-world GitHub issues? In: \emph{International Conference on Learning Representations (ICLR)}.

\bibitem[Anthropic(2026)]{claudeCode}
Anthropic, 2026. Claude Code. Software. \url{https://claude.com/claude-code} (accessed 18 September 2026). %% [VERIFY] preferred citation form

\bibitem[Google(2026)]{jules}
Google, 2026. Jules: an asynchronous coding agent. Software. \url{https://jules.google/} (accessed 18 September 2026).

%% --- Policies and platform documentation ------------------------------------
\bibitem[Elsevier(2026)]{elsevierAI}
Elsevier, 2026. Generative AI policies for journals. \url{https://www.elsevier.com/about/policies-and-standards/generative-ai-policies-for-journals} (accessed 18 September 2026).

\bibitem[GitHub(2026)]{githubOAuth}
GitHub, 2026. Authorizing OAuth apps. GitHub Docs. \url{https://docs.github.com/en/apps/oauth-apps/building-oauth-apps/authorizing-oauth-apps} (accessed 18 September 2026).

%% -------------------------------------------------------------------------
%% [PLACEHOLDER] Still to be supplied:
%%   \bibitem{energese}  Maud, S., in preparation. The companion energese paper,
%%       Ecological Modelling. Cite by DOI once it exists (Zenodo:
%%       10.5281/zenodo.22340275 for the software).
%%   \bibitem{gssk}  GSSK (General Systems Simulation Kernel): repository URL,
%%       version and DOI.
%%   Citations for Stella, Vensim and Insight Maker (Section 2), and any prior
%%   re-implementation of individual Odum minimodels found in the literature.
%% -------------------------------------------------------------------------

\end{thebibliography}

\end{document}
